\chapter{Introducción}
\label{chap:introduccion}

\section{Contextualización}
El patrimonio cultural de un pueblo no reside únicamente en sus monumentos, sino en los saberes que sus miembros practican y transmiten cada día. La Convención de la UNESCO para la Salvaguardia del Patrimonio Cultural Inmaterial reconoce como tal a los usos, representaciones, expresiones, conocimientos y técnicas que las comunidades heredan de generación en generación y recrean de forma constante en función de su entorno \cite{unesco2003convencion}. Esta definición desplaza la atención desde el objeto hacia el portador del conocimiento y, en consecuencia, plantea que la continuidad del patrimonio depende de la continuidad de su transmisión.

La Sierra Sur del Ecuador constituye un territorio en el que esta dimensión inmaterial adquiere una especial densidad. En él habitan comunidades Kichwa portadoras de conocimientos vinculados con la lengua, la memoria oral, la medicina ancestral, la gastronomía, las festividades, los rituales, la música, la danza, la vestimenta, las técnicas artesanales y las formas de organización comunitaria. Estas manifestaciones se articulan en torno a la cosmovisión andina, de modo que no constituyen expresiones aisladas, sino un sistema que relaciona a las comunidades con el territorio, la naturaleza y la vida colectiva. Experiencias recientes en el pueblo Saraguro, como la revitalización intercultural de la arriería o el registro de rituales ancestrales, confirman la vigencia de este patrimonio vivo en la región \cite{quezada2026muleteering,torres2022ancestral}.

No obstante, la vigencia de estos saberes no garantiza su permanencia. Los procesos de migración, la globalización, la transformación de las dinámicas sociales, la disminución del uso cotidiano de la lengua kichwa y la escasa participación de las nuevas generaciones debilitan los canales tradicionales de transmisión intergeneracional. Cuando el conocimiento se conserva principalmente en la memoria de sabios, líderes comunitarios, artesanos, músicos y agricultores, el envejecimiento de estos depositarios se traduce en un riesgo directo de pérdida, descontextualización o transformación de las manifestaciones culturales.

En este escenario, las tecnologías de la información ofrecen alternativas para registrar, organizar y difundir el patrimonio mediante aplicaciones web y móviles, contenidos multimedia, servicios de georreferenciación y recursos educativos interactivos. Sin embargo, su incorporación exige atender las condiciones reales de las comunidades: la conectividad limitada de los sectores rurales, la diversidad lingüística, la accesibilidad y la protección de conocimientos considerados sensibles o de uso restringido \cite{christen2012information}. Es en esta intersección entre patrimonio vivo y tecnología donde se sitúa el presente trabajo.

El proyecto se desarrolla en el marco de la Universidad Tecnológica Indoamérica, a través del Centro de Investigación de Ciencias Humanas y de la Educación (CICHE), como parte del proyecto Kallarimanta. Esta iniciativa se orienta a estudiar, documentar y revitalizar la cultura y las tradiciones de los pueblos Kichwa de la Sierra Sur del Ecuador, articulando la investigación intercultural con la preservación del patrimonio cultural inmaterial y su proyección educativa. La solución propuesta, denominada Kallarimanta Digital, busca integrar los contenidos recopilados por el proyecto dentro de un sistema tecnológico estructurado, seguro y escalable.


\section{Estado del Arte}
La Convención de la UNESCO para la Salvaguardia del Patrimonio Cultural Inmaterial reconoce como patrimonio a los usos, representaciones, expresiones, conocimientos y técnicas que las comunidades transmiten de generación en generación \cite{unesco2003convencion}. A partir de este marco, la revisión de la literatura se organizó en tres ejes: (i) aplicaciones digitales orientadas al patrimonio Kichwa y andino, (ii) gobernanza y soberanía de los datos indígenas, y (iii) aprendizaje móvil y georreferenciación aplicados a la difusión cultural.

\subsection{Soluciones digitales para el patrimonio Kichwa y andino}
El antecedente más directo es el trabajo de Quezada-Sarmiento \textit{et al.} \cite{quezada2026mobile}, quienes diseñaron un prototipo móvil educativo para el patrimonio cultural Saraguro y el aprendizaje introductorio del kichwa. La aplicación, desarrollada con la metodología ágil Mobile-D \cite{abrahamsson2004mobiled} sobre una arquitectura Android (MVVM) con servicios de Firebase, organiza contenidos de gastronomía, rituales y festividades, artesanías, medicina ancestral, relatos y lugares culturales en un mapa, además de expresiones básicas en kichwa con audio. La evaluación exploratoria con doce participantes reportó una buena claridad percibida y una alta valoración de los contenidos locales y georreferenciados. Sin embargo, los propios autores señalan limitaciones: ausencia de acceso sin conexión, soporte multilingüe reducido, falta de un protocolo estandarizado de usabilidad y ausencia de un módulo de administración y validación comunitaria de los contenidos.

En la misma línea, Quezada-Sarmiento \textit{et al.} \cite{quezada2026muleteering} abordan la arriería en Saraguro como patrimonio vivo y destacan el diálogo de saberes como estrategia de salvaguardia, mientras que Torres-Toukoumidis \textit{et al.} \cite{torres2022ancestral} analizan los rituales ancestrales de los Andes ecuatorianos como recurso de turismo comunitario. Ambos estudios confirman el valor de estas manifestaciones, aunque no proponen una plataforma tecnológica para su gestión.

\subsection{Gobernanza y soberanía de los datos indígenas}
La digitalización del patrimonio indígena plantea retos que van más allá de lo técnico. Christen \cite{christen2012information} cuestiona la idea de que toda la información deba ser de libre acceso y propone respetar los protocolos culturales de cada comunidad para determinar quién puede ver o compartir ciertos conocimientos. Kukutai y Taylor \cite{kukutai2016indigenous} y Hudson \textit{et al.} \cite{hudson2023indigenous} desarrollan el concepto de soberanía de los datos indígenas, según el cual los pueblos deben conservar el control sobre la información que se genera acerca de ellos. Linkola-Aikio \textit{et al.} \cite{linkola2025digital} complementan esta perspectiva con prácticas sostenibles para el patrimonio cultural indígena digital. Estos trabajos sustentan la necesidad de incorporar roles, permisos, consentimiento y validación comunitaria en la arquitectura propuesta.

\subsection{Aprendizaje móvil y georreferenciación}
Sarrab \textit{et al.} \cite{sarrab2012mobile} describen el aprendizaje móvil como una modalidad que amplía el acceso a recursos educativos en cualquier momento y lugar, lo que resulta pertinente para estudiantes y comunidades dispersas. Por su parte, Gretzel \textit{et al.} \cite{gretzel2015smart} muestran cómo las tecnologías basadas en la ubicación permiten interpretar territorios y recorridos, lo cual respalda el uso de mapas interactivos para vincular cada manifestación cultural con su espacio de origen.

\subsection{Análisis comparativo}
La Tabla~\ref{tab:estado_arte} resume los trabajos revisados, su aporte principal y su relación con el presente proyecto.

\needspace{16\baselineskip}
{\centering\singlespacing
  \captionof{table}{Comparación de los trabajos revisados en el estado del arte.}
  \label{tab:estado_arte}
  \begin{tabular}{|p{0.16\textwidth}|p{0.36\textwidth}|p{0.36\textwidth}|}
    \hline
    \textbf{Trabajo} & \textbf{Aporte principal} & \textbf{Limitación o brecha} \tabularnewline
    \hline
    \raggedright \cite{quezada2026mobile} & \raggedright App móvil con contenidos Saraguro, mapa y kichwa básico & \raggedright Sin modo sin conexión, sin panel web de gestión ni validación comunitaria \tabularnewline
    \hline
    \raggedright \cite{quezada2026muleteering} & \raggedright Revitalización intercultural de la arriería como patrimonio vivo & \raggedright No propone una solución tecnológica \tabularnewline
    \hline
    \raggedright \cite{torres2022ancestral} & \raggedright Rituales ancestrales andinos como turismo comunitario & \raggedright Enfoque turístico, sin plataforma digital \tabularnewline
    \hline
    \raggedright \cite{christen2012information}, \cite{hudson2023indigenous} & \raggedright Protocolos culturales y soberanía de datos indígenas & \raggedright Marco conceptual, sin implementación \tabularnewline
    \hline
    \raggedright \cite{sarrab2012mobile}, \cite{gretzel2015smart} & \raggedright Aprendizaje móvil y tecnologías basadas en ubicación & \raggedright No abordan el patrimonio Kichwa \tabularnewline
    \hline
  \end{tabular}\par}
\elaboradopor{Michael Jhosue Almeida Díaz}
\fuente{Elaboración propia}

En síntesis, las investigaciones revisadas validan la pertinencia de las tecnologías móviles y georreferenciadas para difundir el patrimonio Kichwa, pero ninguna integra en una misma solución una aplicación web de administración, una aplicación móvil con consulta sin conexión, un módulo de georreferenciación y mecanismos de validación comunitaria y control de acceso a contenidos sensibles. El presente proyecto toma como punto de partida el prototipo Saraguro \cite{quezada2026mobile} y busca cubrir esta brecha mediante la arquitectura multiplataforma propuesta.

\section{Planteamiento del Problema}
Si el patrimonio cultural inmaterial existe en la medida en que se transmite, cabe preguntarse: ¿puede preservarse y difundirse de manera efectiva un conocimiento que se encuentra disperso, registrado de forma informal y ausente de los entornos digitales en los que hoy aprenden las nuevas generaciones? Esta pregunta orienta el análisis del problema que motiva el presente trabajo.

La primera dimensión del problema es la dispersión de la información. Una parte importante del conocimiento sobre el patrimonio Kichwa de la Sierra Sur se encuentra distribuida en investigaciones académicas, archivos audiovisuales, documentos institucionales, registros comunitarios y testimonios orales. Al no estar integrados en un sistema común, estos recursos resultan difíciles de localizar, consultar, actualizar y relacionar entre sí, tanto para estudiantes, docentes e investigadores como para los propios miembros de las comunidades.

La segunda dimensión corresponde al carácter informal del registro. Gran parte de los saberes permanece conservada únicamente en la memoria de sabios, líderes comunitarios, artesanos y otros portadores de conocimiento. Esta forma de conservación, legítima dentro de la tradición oral, se vuelve vulnerable cuando coincide con el envejecimiento de los depositarios y la migración de los jóvenes, ya que cada conocimiento no documentado ni transmitido representa una pérdida potencialmente irreversible.

La tercera dimensión es de carácter tecnológico. Existe una limitada disponibilidad de herramientas que permitan recopilar, organizar, validar, georreferenciar y difundir estos conocimientos de forma estructurada e intercultural. Las iniciativas existentes suelen resolver aspectos parciales: el prototipo móvil desarrollado para el patrimonio Saraguro demostró la valoración de los contenidos locales y georreferenciados, pero carece de consulta sin conexión, de soporte multilingüe amplio y de un módulo de administración y validación comunitaria de los contenidos \cite{quezada2026mobile}. A ello se suma que la gestión de información indígena requiere mecanismos de control de acceso y consentimiento que respeten la soberanía de las comunidades sobre sus datos \cite{kukutai2016indigenous,hudson2023indigenous}, aspecto que rara vez se incorpora en el diseño de este tipo de soluciones.

La cuarta dimensión es educativa. Aunque las tecnologías web y móviles amplían las posibilidades de aprendizaje en cualquier momento y lugar \cite{sarrab2012mobile}, todavía son escasos los recursos educativos digitales, bilingües y accesibles que representen adecuadamente los conocimientos de los pueblos Kichwa de la Sierra Sur. En consecuencia, el patrimonio permanece ausente de los entornos en los que estudiantes y jóvenes construyen hoy su aprendizaje.

En síntesis, el problema no radica en la falta de patrimonio, sino en la ausencia de una infraestructura tecnológica integrada, intercultural y culturalmente responsable que permita gestionar, ubicar territorialmente y difundir con fines educativos los saberes Kichwa de la Sierra Sur del Ecuador. Esta carencia mantiene la dispersión de la información, incrementa el riesgo de pérdida del conocimiento y limita su transmisión a las nuevas generaciones.

\section{Justificación}
Frente a la problemática descrita, la pregunta ya no es si la tecnología debe intervenir en la preservación del patrimonio, sino de qué manera puede hacerlo sin sustituir los procesos comunitarios de transmisión ni comprometer los derechos colectivos de sus portadores. El presente proyecto se justifica en la medida en que propone una respuesta técnica a esta cuestión, sustentada en criterios tecnológicos, educativos, interculturales e institucionales.

Desde el ámbito tecnológico, el proyecto integra en una misma arquitectura multiplataforma una aplicación web administrativa, una aplicación móvil, una API de servicios, una base de datos espacial y un repositorio multimedia. La aplicación web se desarrollará con React y TypeScript, la aplicación móvil con Flutter y Dart, y el backend con Node.js, NestJS y TypeScript, organizado como un monolito modular que expone servicios REST documentados mediante OpenAPI/Swagger. La información se almacenará en PostgreSQL con la extensión PostGIS, lo que permitirá registrar coordenadas, representar comunidades y manifestaciones en mapas interactivos con MapLibre y realizar consultas espaciales. A su vez, el uso de SQLite en el dispositivo móvil habilitará la consulta sin conexión de contenidos previamente descargados, respondiendo a la conectividad limitada de los sectores rurales. Esta combinación tecnológica facilita el mantenimiento, la escalabilidad y la incorporación futura de nuevas comunidades, categorías culturales e idiomas.

Desde la perspectiva de la seguridad y la ética de la información, la arquitectura incorpora autenticación mediante JSON Web Token (JWT), control de acceso basado en roles y un flujo de revisión y validación comunitaria previo a la publicación de los contenidos. Estos mecanismos permiten aplicar los principios de protocolos culturales y soberanía de los datos indígenas \cite{christen2012information,hudson2023indigenous}, de modo que la información considerada sensible, reservada o sagrada no se publique sin la autorización de sus portadores. Asimismo, la georreferenciación se gestionará con precaución, evitando exponer la ubicación exacta de lugares ceremoniales o de acceso restringido.

En el ámbito educativo e intercultural, la plataforma transforma los conocimientos culturales en recursos didácticos interactivos, como fichas informativas, contenidos multimedia, mapas y cuestionarios, que podrán utilizarse en procesos de enseñanza relacionados con la identidad, la interculturalidad, la historia local y el patrimonio. El acceso desde dispositivos móviles amplía el alcance de estos recursos hacia estudiantes y jóvenes \cite{sarrab2012mobile}, mientras que la representación territorial de las manifestaciones vincula cada saber con su espacio de origen \cite{gretzel2015smart}. De esta manera, la tecnología no se limita a almacenar información, sino que promueve su comprensión y valoración.

En el plano institucional, el proyecto aporta a la Universidad Tecnológica Indoamérica y al Centro de Investigación de Ciencias Humanas y de la Educación un producto tecnológico aplicable del proyecto Kallarimanta, que podrá emplearse como repositorio, herramienta educativa y medio de difusión de los resultados de investigación. Además, fortalece los procesos de investigación, innovación y vinculación con la sociedad de la institución.

Los beneficiarios directos de la solución son los docentes, investigadores, estudiantes y representantes de las comunidades Kichwa vinculados con el proyecto Kallarimanta, quienes contarán con una herramienta para gestionar y consultar los contenidos. Los beneficiarios indirectos son las nuevas generaciones de las comunidades, el público interesado en la cultura Kichwa y futuros equipos de desarrollo que podrán ampliar la plataforma. En cuanto a la factibilidad, el proyecto es viable desde el punto de vista técnico, al basarse en tecnologías de código abierto ampliamente documentadas, y desde el punto de vista operativo, al contar con el respaldo institucional del CICHE y con el acceso a los materiales y actores del proyecto Kallarimanta.

Finalmente, el proyecto posee pertinencia académica para la carrera de Ingeniería en Tecnologías de la Información, ya que integra el análisis de requerimientos, el diseño de arquitectura de software, el desarrollo web y móvil, la administración de bases de datos espaciales, la seguridad de la información y la evaluación de usabilidad en la resolución de una necesidad real del entorno social.

\section{Objetivos}
\subsection{Objetivo General}
Desarrollar una arquitectura multiplataforma web y móvil que permita gestionar, georreferenciar y difundir con fines educativos el patrimonio cultural inmaterial de los pueblos Kichwa de la Sierra Sur del Ecuador, mediante la integración de contenidos multimedia, recursos interactivos y mecanismos de participación comunitaria, con el propósito de contribuir a la preservación digital, valoración cultural y transmisión intergeneracional de sus saberes y tradiciones.

\subsection{Objetivos Específicos}
\begin{enumerate}
  \item Identificar y analizar los requerimientos funcionales, no funcionales, educativos y culturales de la plataforma, mediante la participación de docentes, investigadores, estudiantes y representantes de las comunidades Kichwa vinculadas con el proyecto Kallarimanta.
  \item Diseñar la arquitectura tecnológica, el modelo de datos y las interfaces de usuario de la solución web y móvil, considerando criterios de modularidad, escalabilidad, accesibilidad, seguridad, pertinencia cultural y conectividad limitada.
  \item Desarrollar una aplicación web administrativa que permita gestionar usuarios, comunidades, categorías, manifestaciones culturales, contenidos multimedia, ubicaciones geográficas y procesos de revisión y validación de la información.
  \item Implementar una aplicación móvil multiplataforma que facilite la consulta de contenidos culturales y educativos, la visualización de mapas interactivos, la búsqueda de información y el acceso limitado sin conexión a Internet.
  \item Integrar un módulo de georreferenciación para registrar y representar territorialmente comunidades, lugares y manifestaciones del patrimonio cultural inmaterial Kichwa de la Sierra Sur del Ecuador.
\end{enumerate}

\section{Alcance del Proyecto}
El proyecto comprende el análisis, diseño, desarrollo y evaluación de un prototipo funcional de la plataforma Kallarimanta Digital, una arquitectura multiplataforma orientada a la gestión, georreferenciación y difusión educativa del patrimonio cultural inmaterial de los pueblos Kichwa de la Sierra Sur del Ecuador. La primera versión se aplicará a una muestra de comunidades, manifestaciones y contenidos vinculados con el proyecto Kallarimanta, y su arquitectura se diseñará con criterios de escalabilidad para incorporar posteriormente nuevas comunidades, categorías culturales, recursos educativos e idiomas.

La solución incluye los siguientes componentes:

\begin{itemize}
  \item \textbf{Aplicación web administrativa:} gestión de usuarios, roles y permisos; registro de comunidades y ubicaciones; administración de categorías culturales; registro, edición, clasificación, revisión, validación y publicación de manifestaciones patrimoniales; carga de fotografías, audios, videos y documentos; restricción del acceso a información cultural sensible; y consulta de reportes básicos.
  \item \textbf{Aplicación móvil multiplataforma:} orientada principalmente a dispositivos Android, permitirá consultar contenidos culturales, reproducir recursos multimedia, explorar comunidades y manifestaciones en un mapa interactivo, buscar información por categorías y palabras clave, acceder a recursos educativos y consultar sin conexión los contenidos previamente descargados.
  \item \textbf{API de servicios:} interfaz REST documentada que centralizará las reglas de negocio, la autenticación, el control de acceso por roles y la comunicación entre las aplicaciones y la base de datos.
  \item \textbf{Base de datos cultural georreferenciada:} almacenamiento estructurado de la información cultural, educativa, geográfica y multimedia, con soporte para contenidos en español y kichwa.
  \item \textbf{Módulo de georreferenciación:} registro y representación territorial de comunidades, lugares y manifestaciones, con tratamiento diferenciado de las ubicaciones consideradas sensibles.
  \item \textbf{Recursos educativos digitales:} fichas informativas, contenidos multimedia y cuestionarios vinculados con las manifestaciones registradas.
  \item \textbf{Repositorio inicial de contenidos:} conjunto de manifestaciones con su descripción, comunidad de procedencia, categoría, ubicación, recursos multimedia, autoría, fuente, condiciones de acceso y estado de validación.
  \item \textbf{Documentación técnica y funcional:} documento de arquitectura, modelo de base de datos, diagramas UML, documentación de la API, código fuente, manual de instalación, manual técnico, manual de usuario e informe de pruebas.
\end{itemize}

La validación del prototipo abarcará pruebas funcionales, de integración, compatibilidad, rendimiento y seguridad, así como una evaluación de usabilidad y aceptación con usuarios representativos, entre ellos estudiantes, docentes, investigadores y miembros de las comunidades participantes.

Quedan fuera del alcance del proyecto los siguientes aspectos:

\begin{itemize}
  \item El desarrollo de versiones nativas independientes para cada sistema operativo. La aplicación móvil se construirá con una tecnología multiplataforma y su funcionamiento en iOS dependerá de la compatibilidad alcanzada.
  \item La digitalización exhaustiva de la totalidad del patrimonio cultural inmaterial de la Sierra Sur del Ecuador.
  \item La sustitución de los procesos comunitarios tradicionales de transmisión de conocimientos.
  \item La publicación de información cultural considerada reservada, sagrada o sensible sin la autorización de sus portadores.
  \item La traducción automática de contenidos al kichwa. La incorporación de textos en esta lengua dependerá de la disponibilidad de especialistas y representantes comunitarios que los traduzcan y validen.
  \item La creación o edición de contenidos desde la aplicación móvil, que en esta versión será exclusivamente de consulta.
  \item La implementación de una arquitectura de microservicios, realidad aumentada, analítica avanzada de uso o mecanismos de edición colaborativa abierta, que se plantean como trabajo futuro.
  \item El alojamiento, mantenimiento y actualización de la plataforma posteriores a la entrega del prototipo, cuya continuidad dependerá del compromiso institucional.
\end{itemize}

\section{Fundamentación Teórica}

\subsection{Patrimonio cultural inmaterial}
La Convención para la Salvaguardia del Patrimonio Cultural Inmaterial define este patrimonio como los usos, representaciones, expresiones, conocimientos y técnicas que las comunidades reconocen como parte de su identidad y que se transmiten y recrean de generación en generación \cite{unesco2003convencion}. La Convención agrupa sus manifestaciones en cinco ámbitos: tradiciones y expresiones orales, incluido el idioma; artes del espectáculo; usos sociales, rituales y actos festivos; conocimientos y usos relacionados con la naturaleza y el universo; y técnicas artesanales tradicionales.

El mismo instrumento entiende la salvaguardia como el conjunto de medidas orientadas a garantizar la viabilidad de este patrimonio, entre ellas su identificación, documentación, investigación, preservación, protección, promoción, valorización y transmisión \cite{unesco2003convencion}. Estas medidas se corresponden con las funciones que la plataforma propuesta debe soportar: registrar, organizar, validar y difundir las manifestaciones culturales sin desvincularlas de sus portadores.

\subsection{Pueblos Kichwa de la Sierra Sur y patrimonio vivo}
En la Sierra Sur del Ecuador, el patrimonio inmaterial Kichwa se expresa en la lengua, la medicina ancestral, la gastronomía, las festividades como el Inti Raymi, las artesanías, los relatos y las formas de organización comunitaria, articuladas en torno a la cosmovisión andina \cite{quezada2026mobile}. Estas prácticas constituyen un patrimonio vivo, cuya continuidad depende del diálogo de saberes entre generaciones y de la participación activa de las comunidades \cite{quezada2026muleteering}. Por esta razón, una solución tecnológica en este ámbito debe concebirse como un apoyo a la transmisión comunitaria y no como su reemplazo.

\subsection{Preservación digital y gobernanza de datos indígenas}
La preservación digital consiste en el conjunto de procesos que permiten registrar, almacenar y mantener accesible la información a lo largo del tiempo. Cuando esta información corresponde a conocimientos indígenas, su gestión plantea exigencias adicionales. Christen \cite{christen2012information} sostiene que el acceso abierto no siempre es adecuado y que deben respetarse los protocolos culturales que determinan quién puede ver o compartir determinados conocimientos.

En esta línea, el concepto de soberanía de los datos indígenas reconoce el derecho de los pueblos a gobernar la información que se genera sobre ellos \cite{kukutai2016indigenous,hudson2023indigenous}. Los principios CARE (beneficio colectivo, autoridad para controlar, responsabilidad y ética) complementan esta perspectiva y orientan la gestión de datos indígenas hacia el bienestar de las comunidades \cite{carroll2020care}. En la plataforma, estos principios se traducen en niveles de acceso, validación comunitaria previa a la publicación y registro del consentimiento y la autoría de cada recurso.

\subsection{Aprendizaje móvil}
El aprendizaje móvil (\textit{m-learning}) es una modalidad educativa que aprovecha los dispositivos portátiles para acceder a contenidos y actividades de aprendizaje en cualquier momento y lugar \cite{sarrab2012mobile}. Su principal ventaja para este proyecto es la posibilidad de llevar recursos educativos sobre el patrimonio a estudiantes y comunidades dispersas geográficamente. En contextos rurales con conectividad limitada, esta ventaja requiere complementarse con un enfoque de funcionamiento sin conexión, en el que la aplicación almacena localmente los contenidos y los sincroniza cuando dispone de acceso a Internet.

\subsection{Sistemas de información geográfica y georreferenciación}
Un sistema de información geográfica (SIG) permite capturar, almacenar, analizar y representar datos asociados a una ubicación en la superficie terrestre \cite{longley2015gis}. La georreferenciación consiste en asignar coordenadas a un objeto dentro de un sistema de referencia, comúnmente el sistema geodésico WGS84, de modo que pueda ubicarse en un mapa. En el modelo vectorial, las entidades se representan como puntos, líneas o polígonos; en este proyecto, las comunidades y lugares se modelan como puntos o polígonos asociados a las manifestaciones culturales.

Para el intercambio de esta información entre el servidor y las aplicaciones se emplea GeoJSON, un formato estándar basado en JSON para representar entidades geográficas \cite{butler2016geojson}. La representación territorial del patrimonio permite relacionar cada manifestación con su espacio de origen y facilita la interpretación de lugares y recorridos culturales \cite{gretzel2015smart}.

\subsection{Arquitectura de software}
La arquitectura de software es el conjunto de estructuras de un sistema, compuestas por sus elementos, las relaciones entre ellos y sus propiedades \cite{bass2021software}. Su definición temprana permite razonar sobre atributos de calidad como la mantenibilidad, la escalabilidad y la seguridad antes de la implementación.

El proyecto adopta una arquitectura cliente-servidor, en la que las aplicaciones web y móvil actúan como clientes de un servidor central. En el servidor se aplica el estilo de monolito modular: una sola aplicación desplegable cuyas funcionalidades se organizan en módulos independientes, como usuarios, comunidades, manifestaciones, multimedia o georreferenciación. Este estilo reduce la complejidad de despliegue frente a una arquitectura de microservicios y conserva una separación clara de responsabilidades.

La comunicación entre clientes y servidor sigue el estilo arquitectónico REST (\textit{Representational State Transfer}), que organiza el sistema en recursos identificados por URL y manipulados mediante operaciones estándar del protocolo HTTP, con interacciones sin estado \cite{fielding2000rest}. Los servicios se documentan con la especificación OpenAPI, lo que facilita su prueba y su integración con otros sistemas.

\subsection{Tecnologías de la solución}
La Tabla~\ref{tab:tecnologias} resume las tecnologías seleccionadas para cada componente de la arquitectura.

\needspace{22\baselineskip}
{\centering\singlespacing
  \captionof{table}{Tecnologías seleccionadas para la solución.}
  \label{tab:tecnologias}
  \begin{tabular}{|p{0.27\textwidth}|p{0.27\textwidth}|p{0.34\textwidth}|}
    \hline
    \textbf{Componente} & \textbf{Tecnología} & \textbf{Uso en el proyecto} \tabularnewline
    \hline
    \raggedright Aplicación web & \raggedright React y TypeScript \cite{react2026docs} & \raggedright Panel administrativo y portal de consulta \tabularnewline
    \hline
    \raggedright Aplicación móvil & \raggedright Flutter y Dart \cite{flutter2026docs} & \raggedright Aplicación multiplataforma orientada a Android \tabularnewline
    \hline
    \raggedright Backend & \raggedright Node.js, NestJS y TypeScript \cite{nestjs2026docs} & \raggedright API REST, reglas de negocio y control de acceso \tabularnewline
    \hline
    \raggedright Base de datos & \raggedright PostgreSQL y PostGIS \cite{postgis2026docs} & \raggedright Datos estructurados y consultas espaciales \tabularnewline
    \hline
    \raggedright Mapas & \raggedright MapLibre \cite{maplibre2026docs} & \raggedright Visualización de comunidades y manifestaciones \tabularnewline
    \hline
    \raggedright Almacenamiento local & \raggedright SQLite & \raggedright Consulta móvil sin conexión \tabularnewline
    \hline
    \raggedright Multimedia & \raggedright Almacenamiento de objetos compatible con S3 o Firebase Storage & \raggedright Fotografías, audios, videos y documentos \tabularnewline
    \hline
    \raggedright Autenticación & \raggedright JSON Web Token (JWT) & \raggedright Inicio de sesión y protección de servicios \tabularnewline
    \hline
    \raggedright Despliegue & \raggedright Docker, Docker Compose y GitHub Actions & \raggedright Entornos reproducibles e integración continua \tabularnewline
    \hline
  \end{tabular}\par}
\elaboradopor{Michael Jhosue Almeida Díaz}
\fuente{Elaboración propia}

React es una biblioteca de JavaScript para construir interfaces de usuario mediante componentes reutilizables, lo que facilita el desarrollo de formularios, catálogos y paneles de administración \cite{react2026docs}. Flutter es un marco de desarrollo que permite compilar aplicaciones para distintos sistemas operativos a partir de una misma base de código en Dart \cite{flutter2026docs}. NestJS es un marco para Node.js que organiza el código del servidor en módulos, controladores y servicios, en concordancia con el estilo de monolito modular \cite{nestjs2026docs}. PostGIS extiende PostgreSQL con tipos de datos, índices y funciones espaciales, lo que permite almacenar geometrías y realizar consultas de proximidad \cite{postgis2026docs}. Finalmente, MapLibre es una biblioteca de código abierto para mostrar mapas interactivos basados en teselas vectoriales \cite{maplibre2026docs}.

\subsection{Seguridad y control de acceso}
La autenticación se basa en JSON Web Token (JWT), un estándar abierto que representa de forma compacta un conjunto de declaraciones firmadas digitalmente, de modo que el servidor puede verificar la identidad del usuario en cada solicitud sin mantener una sesión \cite{jones2015jwt}. La autorización se gestiona mediante control de acceso basado en roles (RBAC), modelo en el que los permisos se asignan a roles y estos a los usuarios, lo que simplifica la administración de privilegios \cite{sandhu1996rbac}. En la plataforma, este modelo se combina con niveles de acceso por contenido para proteger los conocimientos sensibles.

\subsection{Design Science Research y Scrum}
Design Science Research (DSR) es un paradigma de investigación orientado a crear y evaluar artefactos tecnológicos que resuelven problemas identificados en un contexto real \cite{hevner2004design}. Su ciclo comprende la identificación del problema, el diseño y la construcción del artefacto, y su evaluación rigurosa, lo que se ajusta al propósito de este trabajo.

Para la construcción del software se emplea Scrum, un marco ágil que organiza el trabajo en iteraciones de duración fija denominadas \textit{sprints}, al final de las cuales se obtiene un incremento funcional del producto que se revisa y ajusta \cite{schwaber2020scrum}. Su carácter iterativo permite incorporar de manera progresiva la retroalimentación de docentes, investigadores y representantes comunitarios.

\subsection{Usabilidad y escala SUS}
La norma ISO 9241-11 define la usabilidad como el grado en que un sistema puede ser utilizado por usuarios específicos para alcanzar objetivos concretos con efectividad, eficiencia y satisfacción en un contexto de uso determinado \cite{iso2018usability}. Para medirla se utilizará la escala System Usability Scale (SUS), un cuestionario de diez ítems con respuestas tipo Likert de cinco puntos que produce una puntuación entre 0 y 100 \cite{brooke1996sus}. Diversos estudios han establecido rangos de interpretación de esta puntuación \cite{bangor2008sus}, y se considera que un valor de 68 corresponde al promedio de los sistemas evaluados \cite{sauro2011sus}, por lo que se adopta como referencia mínima para el prototipo.
